% ECONOMICS_PROBLEM: {"id": "J61-503", "title": "Migration, selection and local adjustment", "jel_code": "J61", "source_id": "OP-0028"}
% !TeX program = xelatex
\documentclass[11pt,a4paper]{article}
\usepackage[margin=23mm]{geometry}
\usepackage{fontspec}
\setmainfont{Latin Modern Roman}
\usepackage{amsmath,amssymb,mathtools}
\usepackage{xcolor,enumitem,booktabs,longtable,array,xurl,hyperref}
\definecolor{muted}{HTML}{53636C}
\hypersetup{colorlinks=true,urlcolor=blue,linkcolor=blue}
\setlength{\parindent}{0pt}
\setlength{\parskip}{5pt}
\newcommand{\R}{\mathbb R}
\newcommand{\E}{\mathbb E}
\newcommand{\N}{\mathbb N}
\newcommand{\ind}{\mathbf 1}
\newcommand{\Var}{\operatorname{Var}}
\newcommand{\Cov}{\operatorname{Cov}}
\newcommand{\supp}{\operatorname{supp}}
\DeclareMathOperator*{\argmin}{arg\,min}
\DeclareMathOperator*{\argmax}{arg\,max}
\newcommand{\entrymeta}[3]{{\sffamily\small\color{muted}\textbf{\texttt{#1}}\quad|\quad #2\par #3\par}}
\newcommand{\Problemref}[1]{\texttt{#1}}
\newcommand{\JELref}[2]{\texttt{#2} (JEL \texttt{#1})}
\begin{document}
\section*{Migration, selection and local adjustment}
\textbf{Problem ID:} \texttt{J61-503}\quad \textbf{JEL:} \texttt{J61}\par
% BEGIN ECONOMICS STATEMENT
\label{problem:OP-0028}
\label{atlas:prob:L8}
\noindent\textbf{Original atlas ID:} L8 (entry 28). \textbf{Classification:} Editorial dynamic spatial-equilibrium and cohort-policy problem.

\noindent\textbf{Original atlas question:} What are the long-run wage, employment, fiscal and innovation effects of migration after capital and location respond?

\subsection*{Definitions and assumptions}
Fix a receiving economy, a baseline native cohort $\mathcal N_0$, migrant skill types, locations, and a horizon $H$. Let $\pi$ specify admissions, work rights, and the timing of inflows. A dynamic spatial equilibrium model $M$ includes migrant selection, natives' mobility, capital accumulation, firm entry, housing supply, production substitution, and fiscal rules. Require equilibrium existence and a documented selection rule. Let $y_{it}^M(\pi)$ and $e_{it}^M(\pi)\in\{0,1\}$ be real earnings and employment for native $i$, tracked wherever $i$ lives. Let $B_t^M(\pi)$ denote consolidated tax receipts minus public expenditure under a stated accounting convention, and let $I_t^M(\pi)$ be a defined innovation outcome.
\subsection*{Formal research question}
For each nonempty finite baseline group $g\subseteq\mathcal N_0$, identify or sharply bound
\[
 \Theta_g^M(\pi,\pi_0)=\left(
 \frac1{|g|}\mathbb E_M\left[\sum_{i\in g}\sum_{t=0}^H\beta^t
      [y_{it}^M(\pi)-y_{it}^M(\pi_0)]\right],\;
 (\Delta e_{g,t},\Delta B_t,\Delta I_t)_{t=0}^H\right),
\]
where $\beta\in(0,1]$ is fixed, $\Delta e_{g,t}=\mathbb E_M[|g|^{-1}\sum_{i\in g}(e_{it}^M(\pi)-e_{it}^M(\pi_0))]$, $\Delta B_t=\mathbb E_M[B_t^M(\pi)-B_t^M(\pi_0)]$, and $\Delta I_t=\mathbb E_M[I_t^M(\pi)-I_t^M(\pi_0)]$. Determine how the identified set changes when short-run capital and housing are fixed versus when their transition laws operate. Require predicted adjustment paths to be validated on independently observed migration episodes rather than fitted only to one inflow.
\subsection*{Interpretation and scope}
Local wages among current residents mix causal wage responses with changes in who remains. Tracking the baseline cohort avoids silently replacing movers with new residents. Migrant welfare and origin-country consequences are additional targets, not included in native outcomes by definition. Budget effects depend on age profiles, public-service costs, intergovernmental transfers, and the time horizon; tax payments are transfers in an overall welfare calculation. Immigration shocks require a credible assignment argument. Shift-share instruments additionally require restrictions on the origin shocks or exposure shares, treatment timing, and correlated demand shocks; historical shares are not automatically exogenous. Estimated skill-substitution elasticities generate policy predictions only with the announced production and mobility model. Opposite wage estimates across designs can therefore concern different populations and adjustment margins.
\subsection*{Source audit and status}
[S1] supplies a local historical shock, [S2] a skill-group analysis, and [S3] a model-based wage analysis with imperfect substitution. They contain substantive results rather than one unresolved theorem. The integrated spatial, fiscal, and innovation target is editorial. The bibliography cannot certify that every immigration margin remains unsettled in 2026.

\subsection*{References}
\begin{enumerate}[label=\textnormal{[S\arabic*]},leftmargin=*]
\item David Card (1990). \emph{The Impact of the Mariel Boatlift on the Miami Labor Market}. Industrial and Labor Relations Review 43(2), 245--257. \url{https://www.nber.org/papers/w3069}. \textit{Locator:} Primary publisher, working-paper, or author record: bibliographic information and abstract; full-text theorem verification is not claimed. \textit{Role:} Local immigration-shock evidence; does not determine every migration policy effect.
\item George J. Borjas (2003). \emph{The Labor Demand Curve is Downward Sloping: Reexamining the Impact of Immigration on the Labor Market}. The Quarterly Journal of Economics 118(4), 1335--1374. \url{https://academic.oup.com/qje/article-abstract/118/4/1335/1925108}. \textit{Locator:} Primary publisher, working-paper, or author record: bibliographic information and abstract; full-text theorem verification is not claimed. \textit{Role:} Skill-group approach to immigration wage effects under stated production assumptions.
\item Gianmarco I. P. Ottaviano and Giovanni Peri (2012). \emph{Rethinking the Effect of Immigration on Wages}. Journal of the European Economic Association 10(1), 152--197. \url{https://doi.org/10.1111/j.1542-4774.2011.01052.x}. \textit{Locator:} Primary publisher, working-paper, or author record: bibliographic information and abstract; full-text theorem verification is not claimed. \textit{Role:} Model-based long-run wage analysis incorporating imperfect substitution.
\end{enumerate}

\subsection*{Common conventions: Source mathematical conventions}

Unless an entry states otherwise, agent and firm sets are finite. State and action spaces are measurable, strategies and outcome maps are measurable, probability laws are specified, and expectations exist for the targets under discussion. Infinite-horizon discount factors lie in $(0,1)$ and discounted objectives are finite. Norms, tolerances and complexity input sizes must be specified for computational comparisons. Constants or primitives that appear locally are inputs, not empirical facts asserted by this catalogue.

\textbf{Identification.} A statistical model $m\in\mathcal M$ implies an observable law $P_m$ and target $\theta(m)$. At law $P$, its identified set is
\[
\Theta(P;\mathcal M)=\{\theta(m):m\in\mathcal M,\ P_m=P\}.
\]
Point identification holds when this set is a singleton, or equivalently when $P_{m_1}=P_{m_2}$ implies $\theta(m_1)=\theta(m_2)$ for all compatible models. Sharp bounds mean bounds attained, or approached when the boundary is not attained, by compatible models. Writing this set is a definition, not a solution: the substantive task is to establish informative restrictions and derive or estimate the set. An unrestricted-confounding model may give only trivial bounds.

\textbf{Inference.} Identification, estimability and valid uncertainty quantification are separate questions. Fix $\alpha\in(0,1)$. A confidence procedure $C_n$ over a declared law class $\mathcal P$ has asymptotic uniform coverage at level $1-\alpha$ when
\[
\liminf_{n\to\infty}\inf_{P\in\mathcal P}\Pr_P\{\theta(P)\in C_n\}\geq1-\alpha.
\]
For set-identified targets the coverage event must be defined separately, for example coverage of the full identified set. If an entry uses identification-indexed models instead of uniquely indexed laws, the same coverage convention is applied over those models. Pointwise limits do not establish uniform validity, and asymptotic validity does not guarantee finite-sample precision. Sample sizes, dependence structures and weak-identification sequences are part of each application.

\textbf{Counterfactuals.} $Y_i(a)$ is a potential outcome under a specified intervention, including its timing, implementation and versions. It is not $Y_i$ conditional on voluntarily choosing $a$. A full intervention vector permits interference; shorthand scalar treatment notation does not silently impose no interference. Assignment, selection, attrition, anticipation and measurement must be specified. A claim about an instrument's local effect is not automatically a population policy effect.

\textbf{Economic equilibrium.} A counterfactual model includes best responses, constraints, feasibility, market clearing, state transitions and expectations consistent with the stated information. When several equilibria exist, either a declared selector is an input or the target is set-valued. Comparisons across policy regimes must use the stated selection convention. Reduced-form relationships are not presumed invariant after an equilibrium-changing intervention.

\textbf{Optimization and welfare.} Optimization is over the stated feasible set. If attainment is not established, $\sup$ or $\inf$ and $\varepsilon$-optimal choices replace an asserted maximizer or minimizer. For example, with value $V=\sup_{a\in\mathcal A}W(a)<\infty$, an $\varepsilon$-optimal set is $\{a\in\mathcal A:W(a)\geq V-\varepsilon\}$ for $\varepsilon>0$. Benefits, costs, risk and health or environmental outcomes can be aggregated only using declared units and conversion weights. Interpersonal utility weights, the treatment of fiscal transfers and foreign profits, discounting, and the behavioral welfare criterion are explicit inputs. A comparative-static derivative additionally requires existence and appropriate differentiability of the selected counterfactual.

\textbf{Sources.} Local labels [S1], [S2], and so forth restart in every question. Locators identify the relevant abstract, model, section or result at the level actually checked; a bibliographic or abstract-level verification is not represented as inspection of an entire proof. Working papers are distinguished from later publications and changing versions. Source summaries are paraphrased. All URL checks and editorial formulations refer to the audit date stated above.
% END ECONOMICS STATEMENT
\end{document}
